\section{Conclusion}

This paper presented \ours, a MambaVision-based architecture for binary remote-sensing change detection. The method decouples temporal evidence into region-aware and boundary-aware pathways through \drbi, combines signed temporal difference with adaptive receptive-field decoding and CRAM-lite spatial modulation, and applies a bounded residual correction near boundaries. A controlled ablation over seven variants demonstrates that the MambaVision encoder provides the dominant accuracy gain, while signed temporal difference, CRAM-lite modulation, and the combined boundary head plus Sobel-edge loss each contribute additional improvement. On DSIFN-CD (averaged over two runs) the method achieves F1\,=\,95.67 and IoU\,=\,91.71. On WHU-CD, it achieves F1\,=\,95.15 and IoU\,=\,90.76. The results show that \ours{} is competitive with recent Mamba-based change detection methods, including Mamba-CD~\cite{peng2026mambacd}, while providing a different region-boundary temporal refinement design. Future work will extend evaluation to additional binary CD benchmarks, investigate learned boundary gate designs beyond the Sobel prior, and explore applying the region--boundary decomposition to semantic CD settings.
